\documentclass[a4paper,12pt]{article}

\usepackage[T2A]{fontenc}
\usepackage[utf8]{inputenc}
\usepackage[russian]{babel}

\usepackage{amsmath,amssymb,amsthm,mathtools}
\usepackage{helvet}
\renewcommand{\familydefault}{\sfdefault}
\usepackage{icomma}

\usepackage{geometry}
\usepackage{microtype}
\usepackage{tikz}
\usepackage{tkz-euclide}
\usepackage{tabularx,booktabs}
\usepackage{enumitem}
\usepackage{hyperref}
\usepackage{parskip}
\usepackage{fancyhdr}
\usepackage{setspace}
\usepackage{xcolor}

\geometry{margin=2.2cm}
\onehalfspacing

\definecolor{accent}{HTML}{008080}
\definecolor{darkaccent}{HTML}{004D4D}
\definecolor{textgray}{HTML}{333333}
\definecolor{softgray}{HTML}{707070}
\definecolor{lightaccent}{HTML}{DCEEEE}

\hypersetup{
    colorlinks=true,
    linkcolor=darkaccent,
    urlcolor=darkaccent
}

\setlength{\parindent}{0pt}
\setlength{\parskip}{0.6em}

\setlist[itemize]{
    leftmargin=1.5em,
    itemsep=0.25em,
    topsep=0.3em
}
\setlist[enumerate]{
    leftmargin=2em,
    itemsep=0.5em,
    topsep=0.4em
}

\pagestyle{fancy}
\fancyhf{}
\fancyhead[L]{\small\textcolor{softgray}{Николь Саркисьянц}}
\fancyhead[R]{\small\textcolor{softgray}{ЕГЭ по математике}}
\fancyfoot[C]{\small\textcolor{softgray}{\thepage}}
\renewcommand{\headrulewidth}{0.3pt}
\renewcommand{\headrule}{%
  \hbox to\headwidth{\color{lightaccent}\leaders\hrule height \headrulewidth\hfill}%
}

\newcommand{\key}[1]{\textcolor{darkaccent}{\textbf{#1}}}
\newcommand{\formula}[1]{%
    \[
        \displaystyle #1
    \]
}

\begin{document}

% =========================================================
% ТИТУЛЬНЫЙ ЛИСТ
% =========================================================

\begin{titlepage}
\thispagestyle{empty}

\begin{center}

\vspace*{2.2cm}

{\Large\textcolor{accent}{\textbf{ЧЕК-ЛИСТ}}}

\vspace{0.8cm}

{\Huge\bfseries
Повторение алгебры\\[0.3em]
и базовой планиметрии
}

\vspace{0.9cm}

{\Large
Показательные и логарифмические уравнения,\\
уравнения с корнями, площадь треугольника,\\
прямоугольный треугольник,\\
теоремы синусов и косинусов
}

\vspace{1.7cm}

{\large
\textcolor{darkaccent}{
Лёвин Артём Александрович\\
эксклюзивно для Николь Саркисьянц
}
}

\vfill

{\large \today}

\vspace{1.3cm}

\begin{tikzpicture}[x=1cm,y=1cm]
    \draw[
        accent,
        line width=1.2pt
    ]
    (0,0)
    .. controls (3,0.8) and (6,-0.8) ..
    (9,0)
    .. controls (11,0.5) and (13,0.2) ..
    (15,0.55);
\end{tikzpicture}

\end{center}
\end{titlepage}

% =========================================================
% ОСНОВНОЙ ЧЕК-ЛИСТ
% =========================================================

\section*{\textcolor{darkaccent}{1. Показательные уравнения}}

\key{Главная идея:} стремитесь привести обе части уравнения к одинаковому основанию.

Если
\[
a^{f(x)}=a^{g(x)}, \qquad
a>0,\quad a\ne1,
\]
то
\[
f(x)=g(x).
\]

Например,
\[
7^{x-1}=49.
\]
Поскольку
\[
49=7^2,
\]
получаем
\[
x-1=2,
\qquad
x=3.
\]

\subsection*{Как действовать}

\begin{enumerate}
    \item Разложите числа на степени подходящего основания.
    \item Постарайтесь собрать выражения с переменной в степенях.
    \item Приведите обе части к одному основанию.
    \item Приравняйте показатели.
    \item Проверьте арифметику при переносах.
\end{enumerate}

\key{Полезный приём.}
Если встречаются выражения вида
\[
\left(\frac34\right)^{f(x)},
\qquad
\left(\frac43\right)^{g(x)},
\]
помните:
\[
\left(\frac43\right)^t
=
\left(\frac34\right)^{-t}.
\]

Это часто позволяет быстро получить единое основание.

\subsection*{Типичные ошибки}

\begin{itemize}
    \item неверный знак показателя при переходе к обратной дроби;
    \item ошибка при раскрытии скобок в показателе;
    \item перенос слагаемого без смены знака;
    \item попытка ``сократить основания'', которые ещё не приведены к одинаковому виду.
\end{itemize}

% ---------------------------------------------------------

\section*{\textcolor{darkaccent}{2. Уравнения с квадратным корнем}}

Для уравнения
\[
\sqrt{f(x)}=g(x)
\]
корректный равносильный переход имеет вид
\[
\boxed{
\sqrt{f(x)}=g(x)
\iff
\begin{cases}
g(x)\ge0,\\
f(x)=g^2(x).
\end{cases}}
\]

Условие
\[
g(x)\ge0
\]
обязательно: арифметический квадратный корень неотрицателен.

\subsection*{Пример}

Решим
\[
\sqrt{6x+13}=x+1.
\]

Получаем
\[
\begin{cases}
x+1\ge0,\\
6x+13=(x+1)^2.
\end{cases}
\]

Из второго уравнения:
\[
6x+13=x^2+2x+1,
\]
\[
x^2-4x-12=0,
\]
\[
(x-6)(x+2)=0.
\]

Кандидаты:
\[
x=6,\qquad x=-2.
\]

Но
\[
x+1\ge0
\iff
x\ge-1.
\]

Следовательно,
\[
\boxed{x=6}.
\]

\subsection*{Ещё один характерный пример}

\[
\sqrt{x+6}=-x.
\]

Необходимо:
\[
-x\ge0
\iff
x\le0.
\]

После возведения в квадрат:
\[
x+6=x^2,
\]
\[
x^2-x-6=0,
\]
\[
(x-3)(x+2)=0.
\]

Получаем кандидатов
\[
3,\quad -2.
\]

Условие \(x\le0\) проходит только число \(-2\).

Поэтому
\[
\boxed{x=-2}.
\]

\key{Главное правило:} возведение в квадрат может породить посторонние корни. Условия и проверка здесь принципиальны.

% ---------------------------------------------------------

\section*{\textcolor{darkaccent}{3. Логарифмические уравнения}}

Для выражения
\[
\log_a f(x)
\]
должны выполняться условия
\[
\boxed{
a>0,\qquad
a\ne1,\qquad
f(x)>0.}
\]

Если основание является постоянным корректным числом, например \(2\), отдельно исследовать его не требуется.

\subsection*{Основной переход}

\[
\log_a f(x)=b
\iff
f(x)=a^b,
\]
при
\[
a>0,\qquad a\ne1,\qquad f(x)>0.
\]

Например,
\[
\log_2(4-x)=7.
\]

Получаем
\[
4-x=2^7=128,
\]
откуда
\[
x=-124.
\]

Проверка аргумента:
\[
4-(-124)=128>0.
\]

Следовательно,
\[
\boxed{x=-124}.
\]

В этом примере положительность аргумента фактически обеспечивается самим равенством
\[
4-x=2^7>0.
\]

\subsection*{Что контролировать}

\begin{itemize}
    \item аргумент каждого логарифма строго положителен;
    \item основание положительно и отлично от единицы;
    \item при преобразованиях не потеряны ограничения;
    \item после решения полученные значения удовлетворяют исходному уравнению.
\end{itemize}

% =========================================================

\newpage
\section*{\textcolor{darkaccent}{4. Площадь треугольника: пять ключевых формул}}

Эти формулы необходимо знать уверенно: метод равенства площадей позволяет связывать высоты, стороны, углы и радиусы окружностей.

\subsection*{1. Через основание и высоту}

\[
\boxed{S=\frac12 ah_a}
\]

Здесь \(h_a\) --- высота, проведённая к стороне \(a\).

\begin{center}
\begin{tikzpicture}[scale=0.9]
    \tkzSetUpLine[line width=1pt,color=accent]
    \tkzSetUpPoint[color=accent,fill=accent,size=3]
    \tkzSetUpLabel[color=black,font=\small]

    \tkzDefPoint(0,0){A}
    \tkzDefPoint(6,0){C}
    \tkzDefPoint(2.2,3.2){B}
    \tkzDefPoint(2.2,0){H}

    \tkzDrawPolygon(A,B,C)
    \tkzDrawSegment[dashed](B,H)

    \tkzMarkRightAngle[size=0.28](B,H,C)

    \tkzDrawPoints(A,B,C,H)
    \tkzLabelPoints[below](A,H,C)
    \tkzLabelPoints[above](B)

    \tkzLabelSegment[below](A,C){$a$}
    \tkzLabelSegment[left](B,H){$h_a$}
\end{tikzpicture}
\end{center}

% ---------------------------------------------------------

\subsection*{2. Через две стороны и угол между ними}

Если угол \(\gamma\) заключён между сторонами \(a\) и \(b\), то
\[
\boxed{
S=\frac12 ab\sin\gamma
}
\]

Аналогично:
\[
S=\frac12 bc\sin\alpha,
\qquad
S=\frac12 ac\sin\beta.
\]

\key{Важно:} берутся именно две стороны и синус угла \emph{между ними}.

% ---------------------------------------------------------

\subsection*{3. Формула Герона}

Если известны три стороны \(a,b,c\), сначала вычисляется полупериметр:
\[
p=\frac{a+b+c}{2}.
\]

Затем:
\[
\boxed{
S=\sqrt{p(p-a)(p-b)(p-c)}
}
\]

Формула особенно полезна, когда известны стороны, но отсутствуют удобные высоты или углы.

% ---------------------------------------------------------

\subsection*{4. Через радиус вписанной окружности}

Для треугольника:
\[
\boxed{S=pr}
\]

где
\[
p=\frac{a+b+c}{2},
\]
а \(r\) --- радиус вписанной окружности.

Эта идея обобщается на многоугольник, около которого можно описать окружность как вписанную в сам многоугольник: площадь равна произведению радиуса вписанной окружности на полупериметр.

% ---------------------------------------------------------

\subsection*{5. Через радиус описанной окружности}

\[
\boxed{
S=\frac{abc}{4R}
}
\]

где \(R\) --- радиус окружности, описанной около треугольника.

\subsection*{Метод равенства площадей}

Одна и та же площадь может быть выражена несколькими способами.

Поэтому допустимо записать, например,
\[
\frac12 ah_a
=
\frac12 bc\sin\alpha,
\]
или
\[
pr
=
\frac{abc}{4R}.
\]

Именно эта идея часто позволяет найти неизвестную высоту или радиус.

% =========================================================

\newpage
\section*{\textcolor{darkaccent}{5. Прямоугольный треугольник}}

Пусть катеты равны \(a\) и \(b\), гипотенуза --- \(c\).

По теореме Пифагора
\[
\boxed{c^2=a^2+b^2}.
\]

Площадь:
\[
S=\frac12 ab.
\]

Если \(h\) --- высота к гипотенузе, то одновременно
\[
S=\frac12 ch.
\]

Следовательно,
\[
\frac12 ab=\frac12 ch,
\]
откуда
\[
\boxed{
h=\frac{ab}{c}
}.
\]

\subsection*{Пример}

Для треугольника с катетами \(6\) и \(8\):
\[
c=\sqrt{6^2+8^2}=10.
\]

Высота к гипотенузе:
\[
h=\frac{6\cdot8}{10}
=
\frac{24}{5}
=
4{,}8.
\]

% ---------------------------------------------------------

\subsection*{Радиусы окружностей: пример \(3\)--\(4\)--\(5\)}

Для прямоугольного треугольника со сторонами
\[
3,\quad4,\quad5
\]
площадь равна
\[
S=\frac12\cdot3\cdot4=6.
\]

Через описанную окружность:
\[
6=\frac{3\cdot4\cdot5}{4R},
\]
поэтому
\[
R=\frac52.
\]

Через вписанную окружность:
\[
p=\frac{3+4+5}{2}=6,
\]
\[
6=6r,
\]
следовательно,
\[
r=1.
\]

Для прямоугольного треугольника также полезно помнить:
\[
\boxed{R=\frac{c}{2}}.
\]

% ---------------------------------------------------------

\subsection*{Медиана к гипотенузе}

\begin{center}
\begin{tikzpicture}[scale=0.9]
    \tkzSetUpLine[line width=1pt,color=accent]
    \tkzSetUpPoint[color=accent,fill=accent,size=3]
    \tkzSetUpLabel[color=black,font=\small]

    \tkzDefPoint(0,0){C}
    \tkzDefPoint(6,0){B}
    \tkzDefPoint(0,4){A}
    \tkzDefMidPoint(A,B)
    \tkzGetPoint{M}

    \tkzDrawPolygon(A,B,C)
    \tkzDrawSegment(C,M)

    \tkzMarkRightAngle[size=0.3](A,C,B)

    \tkzDrawPoints(A,B,C,M)
    \tkzLabelPoints[left](A,C)
    \tkzLabelPoints[right](B)
    \tkzLabelPoints[above right](M)
\end{tikzpicture}
\end{center}

Если \(CM\) --- медиана, проведённая из вершины прямого угла \(C\) к гипотенузе \(AB\), то
\[
\boxed{
AM=BM=CM=\frac{AB}{2}.
}
\]

Следовательно, треугольники \(ACM\) и \(BCM\) равнобедренные.

Например, если
\[
\angle B=32^\circ,
\]
то в равнобедренном треугольнике \(BCM\)
\[
\angle BCM=\angle CBM=32^\circ.
\]

Поскольку
\[
\angle ACB=90^\circ,
\]
получаем
\[
\angle ACM=90^\circ-32^\circ=58^\circ.
\]

\key{Это одно из важнейших специальных свойств прямоугольного треугольника.}

% =========================================================

\newpage
\section*{\textcolor{darkaccent}{6. Высота, биссектриса и медиана из вершины прямого угла}}

Пусть
\[
\angle C=90^\circ,
\qquad
\angle B=32^\circ.
\]

Тогда
\[
\angle A=58^\circ.
\]

Если \(CK\) --- биссектриса угла \(C\), то
\[
\angle ACK=\angle KCB=45^\circ.
\]

Если \(CH\) --- высота к гипотенузе, то из прямоугольного треугольника \(BCH\):
\[
\angle BCH
=
90^\circ-\angle B
=
58^\circ.
\]

Поэтому угол между биссектрисой и высотой:
\[
\angle KCH
=
58^\circ-45^\circ
=
13^\circ.
\]

Если \(CM\) --- медиана к гипотенузе, то
\[
\angle MCB=\angle B=32^\circ.
\]

Следовательно, угол между медианой и высотой:
\[
\angle MCH
=
58^\circ-32^\circ
=
26^\circ.
\]

\subsection*{Почему рисунок имеет значение}

Чертёж не является доказательством, однако он должен сохранять правильный порядок объектов.

Если при вычислении геометрического угла получается отрицательное число, следует проверить:

\begin{itemize}
    \item верно ли расположены высота, медиана и биссектриса;
    \item соответствует ли рисунок данным углам;
    \item какую величину следует вычитать из какой;
    \item действительно ли изображён нужный угол.
\end{itemize}

Угол около \(30^\circ\) должен визуально быть заметно меньше угла около \(60^\circ\)--\(70^\circ\).

Точный чертёж от руки не требуется, но взаимное расположение элементов должно соответствовать условию.

% =========================================================

\newpage
\section*{\textcolor{darkaccent}{7. Ищите прямоугольный треугольник}}

Очень сильная стратегия в планиметрии:

\begin{center}
\large
\textcolor{darkaccent}{
\textbf{найдите готовый прямоугольный треугольник\\
или создайте его дополнительным построением.}
}
\end{center}

Причина проста: для прямоугольного треугольника доступны

\begin{itemize}
    \item теорема Пифагора;
    \item определения \(\sin\), \(\cos\), \(\tg\);
    \item простая формула площади;
    \item специальное свойство медианы к гипотенузе;
    \item подобие возникающих треугольников.
\end{itemize}

\subsection*{Пример с высотой}

Пусть
\[
AC=14,
\qquad
\sin\angle A=\frac{12}{13},
\]
а \(CH\) --- высота к \(AB\).

В прямоугольном треугольнике \(ACH\)
\[
\sin\angle A
=
\frac{CH}{AC}.
\]

Следовательно,
\[
\frac{12}{13}
=
\frac{CH}{14},
\]
откуда
\[
\boxed{
CH=\frac{168}{13}.
}
\]

Если в условии дополнительно дана сторона \(BC\), она может оказаться избыточной.

\key{Не пытайтесь обязательно использовать все данные задачи.}

% =========================================================

\newpage
\section*{\textcolor{darkaccent}{8. Теорема синусов}}

Для произвольного треугольника \(ABC\), где стороны
\[
a=BC,\qquad b=AC,\qquad c=AB
\]
лежат напротив углов
\[
A,\qquad B,\qquad C,
\]
выполняется
\[
\boxed{
\frac{a}{\sin A}
=
\frac{b}{\sin B}
=
\frac{c}{\sin C}
=
2R.
}
\]

Здесь \(R\) --- радиус описанной окружности.

\subsection*{Когда теорема особенно полезна}

Если известны:

\begin{itemize}
    \item сторона и противолежащий ей угол;
    \item ещё одна сторона;
    \item либо ещё один угол.
\end{itemize}

Теорема связывает в одной формуле стороны, углы и радиус описанной окружности.

\key{Контроль соответствия:} каждая сторона делится на синус именно \emph{противолежащего} угла.

% =========================================================

\section*{\textcolor{darkaccent}{9. Теорема косинусов}}

Для стороны \(a\), лежащей напротив угла \(A\):
\[
\boxed{
a^2=b^2+c^2-2bc\cos A.
}
\]

Аналогично:
\[
b^2=a^2+c^2-2ac\cos B,
\]
\[
c^2=a^2+b^2-2ab\cos C.
\]

\subsection*{Как выбрать нужную формулу}

Если требуется найти угол, записывайте теорему косинусов для стороны, лежащей \textbf{напротив этого угла}.

\subsection*{Пример: стороны \(3\), \(4\), \(5\)}

Пусть требуется найти угол \(A\), заключённый между сторонами \(3\) и \(5\), а противоположная сторона равна \(4\).

Тогда
\[
4^2
=
3^2+5^2-2\cdot3\cdot5\cos A.
\]

Отсюда
\[
16=34-30\cos A,
\]
\[
30\cos A=18,
\]
\[
\boxed{\cos A=\frac35=0{,}6}.
\]

Если нужен сам геометрический угол:
\[
\boxed{
A=\arccos 0{,}6.
}
\]

В геометрии рассматривается угол треугольника
\[
0^\circ<A<180^\circ,
\]
поэтому общий тригонометрический набор решений здесь не записывается.

% =========================================================

\newpage
\section*{\textcolor{darkaccent}{10. Синус по известному косинусу}}

Основное тригонометрическое тождество:
\[
\boxed{
\sin^2\alpha+\cos^2\alpha=1.
}
\]

Отсюда
\[
\sin^2\alpha=1-\cos^2\alpha.
\]

При решении уравнения относительно неизвестного числа:
\[
\sin\alpha
=
\pm\sqrt{1-\cos^2\alpha}.
\]

Но для внутреннего угла треугольника
\[
0^\circ<\alpha<180^\circ,
\]
поэтому
\[
\sin\alpha>0.
\]

Следовательно, в геометрической задаче:
\[
\boxed{
\sin\alpha
=
\sqrt{1-\cos^2\alpha}.
}
\]

Например, если
\[
\cos\alpha=0{,}8,
\]
то
\[
\sin\alpha
=
\sqrt{1-0{,}8^2}
=
\sqrt{1-0{,}64}
=
\sqrt{0{,}36}
=
0{,}6.
\]

\subsection*{Критически важная алгебраическая ошибка}

В общем случае
\[
\boxed{
\sqrt{a^2-b^2}\ne a-b.
}
\]

Квадратный корень нельзя ``распределять'' по сумме или разности:
\[
\sqrt{x+y}\ne\sqrt{x}+\sqrt{y},
\]
\[
\sqrt{x-y}\ne\sqrt{x}-\sqrt{y}.
\]

Сначала вычисляется всё подкоренное выражение, затем извлекается корень.

При этом
\[
\sqrt{a^2}=|a|.
\]

Для произведения неотрицательных чисел допустимо:
\[
\sqrt{uv}=\sqrt{u}\sqrt{v},
\qquad u\ge0,\quad v\ge0.
\]

Также
\[
\sqrt{a^2b^2}=|ab|.
\]

\key{Знак \(\pm\) появляется при решении уравнения}
\[
z^2=A,
\]
поскольку тогда
\[
z=\pm\sqrt A.
\]

Сам символ
\[
\sqrt A
\]
обозначает только неотрицательное значение.

% =========================================================

\section*{\textcolor{darkaccent}{11. Что необходимо знать наизусть}}

\begin{enumerate}
    \item Условия существования логарифма.
    \item Равносильный переход для
    \[
    \sqrt{f(x)}=g(x).
    \]
    \item Пять основных формул площади треугольника.
    \item Формулу
    \[
    h=\frac{ab}{c}
    \]
    для высоты к гипотенузе прямоугольного треугольника.
    \item Свойство медианы к гипотенузе:
    \[
    AM=BM=CM.
    \]
    \item Определения синуса и косинуса в прямоугольном треугольнике.
    \item Теорему синусов.
    \item Теорему косинусов.
    \item Основное тригонометрическое тождество.
    \item Правило:
    \[
    \sqrt{a^2-b^2}\ne a-b.
    \]
\end{enumerate}

\subsection*{Стратегия при решении геометрии}

\begin{enumerate}
    \item Сделайте рисунок, приблизительно соответствующий условию.
    \item Отметьте все известные углы и длины.
    \item Найдите прямоугольные треугольники.
    \item Проверьте специальные свойства: медиана, биссектриса, высота.
    \item Посмотрите на набор данных:
    \begin{itemize}
        \item основание + высота \(\rightarrow\) площадь;
        \item две стороны + угол между ними \(\rightarrow\) площадь;
        \item три стороны \(\rightarrow\) Герон или теорема косинусов;
        \item сторона + противолежащий угол \(\rightarrow\) теорема синусов;
        \item вписанная окружность \(\rightarrow S=pr\);
        \item описанная окружность \(\rightarrow S=\dfrac{abc}{4R}\).
    \end{itemize}
    \item Если нужно найти высоту или радиус, попробуйте \key{метод равенства площадей}.
    \item Проверьте правдоподобие ответа.
\end{enumerate}

% =========================================================
% ТРЕНИРОВОЧНЫЙ БЛОК
% =========================================================

\newpage
\section*{\textcolor{darkaccent}{Тренировочный блок}}

\textbf{Путь А.} Выполните все 10 заданий без просмотра ответов.

\subsection*{Средний уровень}

\begin{enumerate}[label=\textbf{\arabic*.}]

\item Решите уравнение:
\[
9^{x-2}=27.
\]

\item Найдите корень уравнения:
\[
\sqrt{x+12}=6-x.
\]

\item В прямоугольном треугольнике катеты равны \(9\) и \(12\).
Найдите:
\[
\text{а) гипотенузу; \qquad б) высоту к гипотенузе.}
\]

\end{enumerate}

\subsection*{Выше среднего}

\begin{enumerate}[label=\textbf{\arabic*.},resume]

\item Решите уравнение:
\[
\log_3(11-x)=2.
\]

\item В треугольнике \(ABC\)
\[
AC=15,
\qquad
\sin A=\frac45.
\]
Высота \(CH\) проведена к стороне \(AB\).
Найдите \(CH\).

\item В прямоугольном треугольнике \(ABC\)
\[
\angle C=90^\circ,
\qquad
\angle B=36^\circ.
\]
\(CM\) --- медиана к гипотенузе.
Найдите
\[
\angle ACM.
\]

\item В прямоугольном треугольнике \(ABC\)
\[
\angle C=90^\circ,
\qquad
\angle B=34^\circ.
\]
Из вершины \(C\) проведены высота \(CH\) и биссектриса \(CK\).
Найдите острый угол между \(CH\) и \(CK\).

\item Стороны треугольника равны
\[
5,\qquad7,\qquad8.
\]
Найдите косинус угла, лежащего напротив стороны \(8\).

\item В треугольнике
\[
a=6,\qquad
A=30^\circ,\qquad
B=45^\circ.
\]
Найдите сторону \(b\).

\end{enumerate}

\subsection*{Сложное задание}

\begin{enumerate}[label=\textbf{\arabic*.},resume]

\item В прямоугольном треугольнике \(ABC\)
\[
\angle C=90^\circ,
\qquad
AC=8,
\qquad
BC=15.
\]

Из вершины \(C\) проведены высота \(CH\) к гипотенузе и медиана \(CM\).

Найдите:

\[
\text{а) }AB;
\qquad
\text{б) }CH;
\qquad
\text{в) }R;
\qquad
\text{г) }r;
\qquad
\text{д) }\cos\angle MCH.
\]

\end{enumerate}

% =========================================================
% ОТВЕТЫ
% =========================================================

\newpage
\section*{\textcolor{darkaccent}{Ответы к тренировочному блоку}}

\renewcommand{\arraystretch}{1.45}

\begin{table}[ht]
\centering
\caption{Краткие ответы}
\begin{tabularx}{\linewidth}{@{}cX@{}}
\toprule
\textbf{№} & \textbf{Ответ} \\
\midrule

1 &
\[
x=\frac72
\]
\\

2 &
\[
x=4
\]
\\

3 &
\[
\text{а) }15;
\qquad
\text{б) }\frac{36}{5}
\]
\\

4 &
\[
x=2
\]
\\

5 &
\[
CH=12
\]
\\

6 &
\[
\angle ACM=54^\circ
\]
\\

7 &
\[
11^\circ
\]
\\

8 &
\[
\cos\alpha=\frac17
\]
\\

9 &
\[
b=6\sqrt2
\]
\\

10 &
\[
\text{а) }AB=17;
\qquad
\text{б) }CH=\frac{120}{17};
\qquad
\text{в) }R=\frac{17}{2};
\]
\[
\text{г) }r=3;
\qquad
\text{д) }\cos\angle MCH=\frac{161}{289}.
\]
\\

\bottomrule
\end{tabularx}
\end{table}

\end{document}